% ============================================================
% Chapter 2 -- Statistical mechanics, energy-based models, and diffusion
%
% The context chapter. It exists to make one distinction available to the rest
% of the thesis: the exact score, the empirical score, and the fitted score are
% three different objects, and the difficulties of generative diffusion are
% properties of the gaps between them. Chapter 6 onwards adds a fourth.
% ============================================================

\chapter{Statistical Mechanics, Energy-Based Models, and Diffusion}
\label{ch:statmech}

\section{What statistical mechanics is, and why it is here}
\label{sec:sm-what}

Newton's second law is deterministic and complete. Given the positions and
momenta of $N$ particles and the forces between them, the trajectory of the
system is fixed for all time. For $N$ of order $10^{23}$ this completeness is
useless: the initial condition cannot be written down, the equations cannot be
integrated, and the answer, if it were available, would be a list of
$6 \times 10^{23}$ numbers and not an explanation of anything. Yet the gas
in a room has a temperature, a pressure and a heat capacity, and these are
stable, reproducible, and predictable from a handful of parameters.

Statistical mechanics is the resolution of that tension. It abandons the
trajectory and keeps a probability distribution over microstates, and it fixes
that distribution by a variational principle: among all distributions with a
given mean energy, the one realised at equilibrium maximises the entropy
\begin{equation}
  S[p] \;=\; -\sum_{x} p(x) \log p(x),
  \label{eq:sm-entropy}
\end{equation}
which is the unique measure of uncertainty consistent with a short list of
consistency axioms \citep{jaynes1957information}. Carrying out the maximisation
subject to $\sum_x p(x) E(x) = \bar{E}$ produces the Boltzmann distribution
\begin{equation}
  p(x) \;=\; \frac{1}{Z(\beta)}\, e^{-\beta E(x)},
  \qquad
  Z(\beta) \;=\; \sum_{x} e^{-\beta E(x)},
  \label{eq:sm-boltzmann}
\end{equation}
with the inverse temperature $\beta$ appearing as the Lagrange multiplier of
the energy constraint \citep{landau1980statistical, goldstein2002classical}.
Everything thermodynamic follows from the normalising constant: the free energy
is $F = -\beta^{-1}\log Z$, and its derivatives give the energy, the entropy and
the response functions.

Two features of this construction matter for what follows, and neither is
specific to physics. The first is that the whole of the macroscopic behaviour
is encoded in a single scalar function of the parameters, $\log Z$, which is a
sum over exponentially many configurations and is therefore hard to compute.
The second is that competition between energy, which orders, and entropy, which
disorders, can produce \emph{phase transitions}: the equilibrium state can
change character abruptly as a parameter is varied, even though the microscopic
law does not change at all. A system with a symmetry can settle into a state
that lacks it, and which of the equivalent states it settles into is decided by
fluctuations. The canonical example is the Ising model \citep{ising1925beitrag},
whose two-dimensional case \citet{onsager1944crystal} solved exactly, and which
remains the reference picture for a transition driven by the competition between
energy and entropy. Both features will reappear, unchanged in substance, in the
generative models of the next sections.

\section{Energy-based models: the same object, learned}
\label{sec:sm-ebm}

Traffic between statistical mechanics and learning machines runs in both
directions and is old. The Hopfield network stored memories as minima of an
energy \citep{hopfield1982neural}; its capacity was settled by spin-glass methods
\citep{amit1985storing, mezard1987spin, gardner1988space}; and the Boltzmann
machine and its restricted form made the Boltzmann distribution itself the model
\citep{ackley1985learning, smolensky1986information}. The 2024 Nobel Prize in
Physics recognised that lineage \citep{nobel2024physics}. What follows takes
\eqref{eq:sm-boltzmann} and reads it in the opposite direction. Rather than deriving a distribution from a known energy, it
posits a parametric energy $E_\theta$ and fits it so that the induced
distribution matches data \citep{ackley1985learning, lecun2006tutorial}:
\begin{equation}
  p_\theta(x) \;=\; \frac{e^{-E_\theta(x)}}{Z(\theta)},
  \qquad
  Z(\theta) \;=\; \int_{\R^{d}} e^{-E_\theta(x)}\, \dd x .
  \label{eq:sm-ebm}
\end{equation}
The microstate becomes a data point, the energy becomes a learnable scoring
function, and the partition function keeps its role and its difficulty. Fitting
by maximum likelihood gives a gradient with the structure of a competition,
\begin{equation}
  \nabla_\theta\, \E_{x \sim p_{\mathrm{data}}}\!\left[-\log p_\theta(x)\right]
  \;=\;
  \E_{x^{+} \sim p_{\mathrm{data}}}\!\left[\nabla_\theta E_\theta(x^{+})\right]
  \;-\;
  \E_{x^{-} \sim p_{\theta}}\!\left[\nabla_\theta E_\theta(x^{-})\right],
  \label{eq:sm-ebm-grad}
\end{equation}
in which the first term lowers the energy of observed data and the second
raises the energy of the model's own samples. The second expectation is over
$p_\theta$, which cannot be sampled directly, so training requires Markov chain
Monte Carlo, typically Langevin dynamics, whose stationary measure is
$\propto e^{-E_\theta}$ \citep{hinton2002training, robert2004monte}. This is
expensive, and its cost is the reason the family was, for a long time,
outperformed by models that give up the explicit density in exchange for cheap
sampling \citep{goodfellow2014generative, kingma2014auto,
papamakarios2021normalizing}. \citet{carbone2025hitchhiker} surveys how the
energy-based view relates to the rest of the generative family.

The learning dynamics of these models are themselves a statistical mechanics
problem, and one with a structure worth carrying forward. The spectral dynamics
of the weights \citep{decelle2017spectral, mezard2017meanfield} and the
emergence of compositional features \citep{tubiana2017emergence} are both
statistical-mechanics accounts of training. \citet{bachtis2024cascade} show that
training a restricted Boltzmann machine
proceeds through a \emph{cascade of second-order phase transitions}: the model
first encodes a coarse, centre-of-mass feature, and then resolves progressively
finer sub-modes, each emergence accompanied by a diverging susceptibility. The
features are learned in order, from the largest scale down. Section
\ref{sec:sm-transitions} finds the same ordering in a different process, and
Chapter~\ref{ch:model} onwards works in a model where the ordering is
analytic.

The connection to what follows is direct. Differentiating
\eqref{eq:sm-ebm} gives
\begin{equation}
  \nabla_x \log p_\theta(x) \;=\; -\nabla_x E_\theta(x),
  \label{eq:sm-score-energy}
\end{equation}
and the intractable $Z(\theta)$, being independent of $x$, has vanished. The
gradient of the log-density carries all the information needed to \emph{sample}
the model while requiring none of the information needed to \emph{normalise}
it. That observation is the whole of the following section.

\section{Generative diffusion}
\label{sec:sm-diffusion}

A diffusion model destroys structure in a controlled way and then learns to run
the destruction backwards \citep{sohl2015deep, song2019generative,
ho2020denoising, song2021scorebased}. Its deterministic counterparts, neural ODEs and flow matching
\citep{chen2018neural, lipman2023flow}, transport the same measures along a
probability-flow trajectory instead of a stochastic one. The forward process is a stochastic differential equation
\begin{equation}
  \dd X_t \;=\; f(X_t, t)\,\dd t \;+\; g(t)\,\dd W_t ,
  \label{eq:sm-forward}
\end{equation}
started at $X_0 \sim p_0$, the data law, and producing marginals $p_t$ that
approach a simple terminal measure. Throughout this thesis the choice is the
Ornstein--Uhlenbeck channel, $f(x,t) = -x$ and $g = \sqrt{2}$, whose transition
law is available in closed form,
\begin{equation}
  x \;=\; e^{-t} a \;+\; \sqrt{\Deltat}\,\varepsilon,
  \qquad \varepsilon \sim \Nd(0, I),
  \qquad \Deltat = 1 - e^{-2t},
  \label{eq:sm-ou}
\end{equation}
and whose stationary law is the standard Gaussian
\citep{uhlenbeck1930theory, gardiner2009stochastic}. Figure
\ref{fig:sm-corruption} shows what \eqref{eq:sm-ou} does to a correlated
sequence: the amplitude contracts, independent noise accumulates, and by
$t \approx 2$ the trajectory is indistinguishable from white noise.

\begin{figure}[t]\centering
\includegraphics[width=\linewidth]{fig_forward_corruption.pdf}
\caption[The forward channel applied to a correlated sequence]{One clean
stationary AR(1) trajectory with $\alpha = 0.85$, in grey, and the same
trajectory after the Ornstein--Uhlenbeck channel \eqref{eq:sm-ou} at four
diffusion times. The signal contracts by $e^{-t}$ while independent noise of
variance $\Deltat$ accumulates; correlation between neighbouring sites survives
longer than amplitude. Generation runs this figure from right to left, and the
score is the only quantity needed to do so.}
\label{fig:sm-corruption}
\end{figure}

The reversal is a classical result. \citet{anderson1982reverse} shows that the
process
\begin{equation}
  \dd X_t \;=\;
  \Big[ f(X_t, t) - g(t)^2\, \nabla_x \log p_t(X_t) \Big]\,\dd t
  \;+\; g(t)\,\dd \bar{W}_t ,
  \label{eq:sm-reverse}
\end{equation}
run backwards from $T$ with $X_T \sim p_T$, reproduces the forward marginals.
Every ingredient of \eqref{eq:sm-reverse} is known by construction except one.
The drift correction is the \emph{score} $\nabla_x \log p_t$, and supplying it
is the entire modelling effort of a diffusion model. A Fokker--Planck argument
for \eqref{eq:sm-reverse} and the probability-flow ODE it induces are in
\compendium{}; a book-length treatment of the model family is
\citet{lai2025principles}.

By \eqref{eq:sm-score-energy} the score is the negative gradient of an energy,
so a diffusion model is an energy-based model that never forms its partition
function and never runs an inner Monte Carlo loop. What it does instead is
train against a target that is available in closed form, which is the content
of the identity in the next section.

\section{Three scores}
\label{sec:sm-three-scores}\label{sec:tweedie}

The word ``score'' is used for three different objects in this literature, and
almost every difficulty in generative diffusion is a property of the gaps
between them. The distinction is drawn sharply in \citet{biroli2023generative}
and organises what follows.

\paragraph{The exact score.} The score of the true noised law,
$\nabla_x \log p_t$ with $p_t = p_0 * \Nd(0,\Deltat I)$ after contraction. It is
the object \eqref{eq:sm-reverse} calls for. It has a form that will be used on
every page of this thesis: writing $p_t$ as the channel convolution and
differentiating under the integral gives
\begin{equation}
  \boxed{\;
  \nabla_x \log p_t(x)
  \;=\; \frac{e^{-t}\,\E[\,a \,|\, x\,] - x}{\Deltat}
  \quad\Longleftrightarrow\quad
  \E[\,a \,|\, x\,]
  = e^{t}\big( x + \Deltat\, \nabla_x \log p_t(x) \big).
  \;}
  \label{eq:tweedie}
\end{equation}
We call \eqref{eq:tweedie} the \emph{score--posterior identity}; in the
empirical-Bayes literature it is Tweedie's formula
\citep{robbins1956empirical, miyasawa1961empirical, efron2011tweedie,
vincent2011connection}. It holds for any $p_0$ whatever, and it says something
stronger than an analogy: computing a score \emph{is} solving a Bayesian
inference problem. The exact score is unavailable in practice for the plain
reason that $p_0$ is unknown.

For \emph{vector-valued} data $a \in \R^{L}$ noised coordinatewise,
\eqref{eq:tweedie} holds componentwise but with the \emph{full} posterior:
\begin{equation}
  \big( \nabla_x \log p_t(x) \big)_k
  \;=\; \frac{e^{-t}\, \E[\,a_k \,|\, x_0, \dots, x_{L-1}\,] - x_k}{\Deltat}.
  \label{eq:tweedie-joint}
\end{equation}
The conditioning is on \emph{every} coordinate. The channel is independent
across frames, but the prior is not, so the posterior mean of frame $k$, and
with it the $k$-th component of the score, depends on all $L$ observations.
Equation~\eqref{eq:tweedie-joint} is the exact sense in which the joint score
of a trajectory differs from a stack of per-frame scores, and the reason this
thesis studies the former.

\paragraph{The empirical score.} Replace $p_0$ by the empirical measure of a
training set $\{a^{(i)}\}_{i=1}^{N}$, which puts mass $1/N$ on each example.
The convolution is then a finite Gaussian mixture and the score is available in
closed form,
\begin{equation}
  \nabla_x \log p^{\mathrm{emp}}_t(x)
  \;=\; \frac{1}{\Deltat}\left(
  e^{-t} \sum_{i=1}^{N} w_i(x)\, a^{(i)} \;-\; x \right),
  \qquad
  w_i(x) \;=\; \frac{e^{-\lVert x - e^{-t} a^{(i)}\rVert^2 / 2\Deltat}}
                    {\sum_{j} e^{-\lVert x - e^{-t} a^{(j)}\rVert^2 / 2\Deltat}},
  \label{eq:sm-empirical}
\end{equation}
a softmax-weighted average of training points. This is exactly computable and
exactly useless as a generative model: as $t \to 0$ the weights concentrate on
the nearest training example and the reverse dynamics deposit the trajectory on
top of it. Running \eqref{eq:sm-reverse} with \eqref{eq:sm-empirical} returns
the training set. Figure~\ref{fig:sm-three-scores} shows the mechanism in one
dimension.

\paragraph{The fitted score.} A neural network $s_\phi(x,t)$ trained by
denoising score matching \citep{hyvarinen2005estimation, vincent2011connection}
to minimise
\begin{equation}
  \E_{t}\, \E_{a \sim p_{\mathrm{data}}}\, \E_{\varepsilon}
  \left\lVert
    s_\phi\big(e^{-t} a + \sqrt{\Deltat}\,\varepsilon,\; t\big)
    + \frac{\varepsilon}{\sqrt{\Deltat}}
  \right\rVert^{2}.
  \label{eq:sm-dsm}
\end{equation}
The minimiser over all measurable functions is the exact score; the minimiser
over a finite training set is the empirical score. Diffusion models work
because the network reaches neither. It is constrained enough to be unable to
represent \eqref{eq:sm-empirical}, and that failure to fit is what produces
novel samples \citep{kadkhodaie2024generalization, kamb2024analytic}. The
mechanism of generalisation is therefore an approximation-theoretic property of
an architecture, which is a considerably less comfortable foundation than the
variational principle of Section~\ref{sec:sm-what}.

\begin{figure}[t]\centering
\includegraphics[width=\linewidth]{fig_three_scores.pdf}
\caption[The exact and empirical scores at three diffusion times]{The exact
score of a standard Gaussian, which is linear in $x$, against the empirical
score \eqref{eq:sm-empirical} built from $N=12$ samples of the same law; ticks
mark the contracted training points $e^{-t}a^{(i)}$. At small $t$ the empirical
score is a sawtooth that points at whichever training example is nearest, and
reverse dynamics driven by it terminate on one of them. As $t$ grows the
weights spread, the sawtooth flattens, and the two scores agree. The useful
regime for a generative model lies between the two, and its width is what
\citet{biroli2024dynamical} quantify.}
\label{fig:sm-three-scores}
\end{figure}

\section{Transitions along the generation trajectory}
\label{sec:sm-transitions}

Because generation is a dynamical process, its failures are dated. Two
crossover times structure the reverse trajectory, and both are statistical
mechanics results in the sense of Section~\ref{sec:sm-what}: a competition
between an entropic and an energetic term, resolved abruptly.

\paragraph{Speciation.} Run backwards from noise, the trajectory commits to a
broad class of outcomes long before it commits to a specific sample. For data
with a covariance whose top eigenvalue is $\lambda$, the commitment happens at
$t_{\mathrm{S}} \sim \tfrac12 \log \lambda$, when the contracted signal along
that direction becomes comparable with the accumulated noise
\citep{biroli2024dynamical, achilli2026speciation}. In high dimension the top
eigenvalue grows with the dimension, so $t_{\mathrm{S}}$ diverges
logarithmically. When the data are hierarchical there is not one such time but
a \emph{ladder} of them, one per level of the hierarchy, and features are
resolved coarse-to-fine \citep{sclocchi2025phase, sclocchi2025probing,
favero2025compositional, favero2025bigger}:
the same ordering that \citet{bachtis2024cascade} observe in the training of an
energy-based model, appearing here along the generation trajectory instead.
Figure~\ref{fig:sm-cascade} shows the ladder in a Gaussian hierarchical model
where every rung can be predicted in closed form from the covariance spectrum
and then measured.

\begin{figure}[t]\centering
\includegraphics[width=\linewidth]{fig_speciation_cascade.pdf}
\caption[The speciation cascade, predicted and measured]{A Gaussian
hierarchical prior has one covariance eigenvalue per level, and therefore a
ladder of speciation times rather than a single one. \emph{Left:} for each
level, the time predicted from the eigenvalue alone (diamonds) against the
times at which the forward and reverse processes are measured to cross
(circles, squares); the grey bar spans the two measurements. \emph{Right:} the
same six levels as prediction against measurement, with the identity line.
Worst-case relative disagreement across all twelve measurements is $3\%$. The
ordering is coarse-to-fine: the largest eigenvalue resolves first.}
\label{fig:sm-cascade}
\end{figure}

\paragraph{Collapse and memorisation.} Later in the reverse trajectory a second
transition occurs, and it is the one that limits the whole construction.
\citet{biroli2024dynamical} and \citet{bonnaire2025memorize} show that dynamics
driven by the empirical score are captured by a single training example past a
time $t_{\mathrm{C}}$, and that avoiding this requires a training set
exponentially large in the data dimension. Memorisation need not be a collapse
onto one example: \citet{achilli2024losing} identify a geometric form in which
dimensions are lost without any single training point capturing the trajectory,
and \citet{garnierbrun2026biased} show how the resulting generalisation can be
biased. Solvable models of the learning problem itself
\citep{cui2025solvable} complete the picture. The criterion is entropic: collapse
occurs when the excess entropy of the noised law relative to the terminal
measure falls below $\log N$, so the dataset size needed to postpone it grows
as $e^{\,\text{(excess entropy)}}$. Figure~\ref{fig:sm-collapse} measures
$t_{\mathrm{C}}$ against $N$ in the chain model of this thesis and compares it
with that closed-form criterion.

\begin{figure}[t]\centering
\includegraphics[width=0.72\linewidth]{fig_collapse_time.pdf}
\caption[Collapse time against dataset size]{Measured collapse time for
empirical-score dynamics on the stationary AR(1) chain, at three chain lengths
$L$, against the number of training sequences. Solid: the time at which the
generated trajectory's nearest-training-neighbour distance falls below the
distance to a fresh sample. Dotted: the excess-entropy criterion of
\citet{biroli2024dynamical}, evaluated in closed form for this model, plotted
where it has a solution. Postponing collapse costs exponentially many
sequences, and the exponent is the excess entropy of the clean law.}
\label{fig:sm-collapse}
\end{figure}

The two transitions bracket the useful part of the trajectory. Above
$t_{\mathrm{S}}$ nothing has been decided; below $t_{\mathrm{C}}$ the empirical
score reproduces its training data. A generative model earns its keep in
between, and the fitted score of \eqref{eq:sm-dsm} is the only one of the three
that operates there, for reasons, as noted above, that are properties of an
architecture rather than of a principle.

\section{Locality, architecture, and message passing}
\label{sec:sm-unet}

If the fitted score is doing the work, the natural question is what the
architecture contributes, and here a second literature meets the first. The
network in almost every image or video diffusion model is a U-Net
\citep{ronneberger2015unet}: a convolutional encoder that repeatedly
downsamples, a decoder that upsamples back, and long skip connections joining
each encoder stage to the decoder stage of the same resolution. The design is
almost universal and was, for a long time, justified anecdotally.

\citet{mei2024unets} gives it a precise reading, and it is the reading this
thesis works inside. For data generated by a tree-structured hierarchical
model, the exact denoiser is computed by belief propagation, which consists of
an upward pass that accumulates evidence from the leaves and a downward pass
that redistributes it, the second reusing the messages formed by the first. A
U-Net has exactly this shape: the encoder is the upward pass, the decoder is
the downward pass, and the long skip connections are the reuse of the
intermediate messages. On this account a U-Net is not merely a good inductive
bias for denoising but an \emph{approximate implementation of the exact
inference algorithm} for hierarchically structured data, which yields sample
complexity bounds for learning the denoiser. Convolutional classifiers admit
the corresponding reading for the upward pass alone.

Two consequences frame this thesis. First, if the architecture is an inference
algorithm, then running the inference algorithm directly is a legitimate
alternative and not a toy substitute, and where the graph is a chain rather
than a tree of higher degree, that algorithm is exact and cheap. Second, the
question of \emph{how large a receptive field a score network needs} becomes
quantitative: it is set by how fast the influence of distant evidence decays in
the corresponding inference problem. Chapter~\ref{ch:gaussian-bp} answers that
question in closed form for the chain. A related result on sequences is
\citet{garnierbrun2024transformers}, who find that a transformer trained on
hierarchically generated sequences reconstructs correlations sequentially,
shortest length scale first, and who use exact belief propagation as the
reference against which the network is measured.

\section{Where this thesis sits}
\label{sec:sm-position}\label{sec:related}\label{sec:gap}

The three scores of Section~\ref{sec:sm-three-scores} exhaust the options
usually considered: one is unavailable, one memorises, and one works for
reasons that are hard to state. This thesis studies a fourth.

For data with sequential structure, a trajectory of a Markov process noised
frame by frame, the score--posterior identity \eqref{eq:tweedie} turns the
score into a posterior expectation on a chain-structured factor graph, and that
posterior is computable exactly by belief propagation. Call the result the
\emph{BP score}. It differs from the other three in a way that is structural
rather than a matter of degree:

\begin{itemize}
  \item Unlike the exact score, it is computable, because the model class is
  specified and the graph is a tree.
  \item Unlike the empirical score, its sufficient statistic is a transition
  kernel, two numbers for a Gaussian AR(1) chain or a fixed grid for a
  nonparametric one, and \emph{not} the training set. The data enter only
  through an average. There is no training point for a trajectory to collapse
  onto, so the memorisation axis of Section~\ref{sec:sm-transitions} does not
  exist for it at any $N$.
  \item Unlike the fitted score, it is transparent: every quantity in it has a
  closed form or a validated numerical reference, so its errors can be
  attributed rather than observed.
\end{itemize}

\noindent The price is an assumption. The BP score is only correct if the data
really are Markov, and the value of the assumption is exactly its correctness:
a claim that has to be measured, not asserted. Chapter~\ref{ch:ring} begins
where the research began, with a two-dimensional dynamic object whose behaviour
motivated studying the joint score at all; Chapters~\ref{ch:gaussian}
to~\ref{ch:laplace} compute the BP score with the transition kernel known;
Chapters~\ref{ch:em-method} and~\ref{ch:em-kernel} learn the kernel from noised
sequences alone; and Chapter~\ref{ch:em-results} measures what the assumption
is worth against fitted scores trained on identical data, and locates where it
stops being worth anything.
